\scanpage{39}
% L1
\uline{gradué} \inserted{\uline{à degrés} $\geq 0$}, et $\xi$ est de degré $1$. De cette façon, $M(X)$ devient gradué [à degrés des deux signes, $\xi$ étant un élément \uline{inversible} \unsure{en|de} degré $1$]. Noter que
\[
M_0(X)=\mathbb Z^{\Sigma_0(X)}
\]
avec
\[
\Sigma_0(X)=\underset{\text{par }x\mapsto x\cdot\xi}{\underrightarrow{\lim}}\Sigma_i^+(X)
\]
et
\[
\begin{aligned}
M^{\mathrm{pair}}(X)&=\underbrace{M_0(X)}_{\text{de degré }0}[\underbrace{\xi}_{\text{de degré }1}]\\
M^{\mathrm{impair}}(X)&=M_1(X)\otimes_{M_0(X)}M_0(X)[\xi]
\end{aligned}
\]
% L10
donc la structure d'anneau gradué de $M(X)$ est déterminée en termes de la structure d'anneau de $M_0(X)$, \deleted{et} le Module $M_1(X)$ sur $M_0(X)$, et l'application de mult.
$M_1(X)\otimes_{M_{0}(X)}M_1(X)\longrightarrow M_0(X)$ \unsure{définie par} \deleted{\unsure{$\xi$}} $u,v\mapsto\xi^{-1}uv$.

% L15
D'autre part, les objets de $D_i^+(M^+(X))$ sont \deleted{\unsure{précist}} \inserted{ceux} dont la classe dans \deleted{$\unsure{M|K}^+(X)$ sont} $M^+(X)$ sont dans $\mathbb Z^{D_i(\Sigma^+(X))}$, où
\[
\left.\begin{aligned}
D_i(\Sigma^+(X))&=\coprod_{j\geq0}\deleted{\unsure{\mathbb Z^{x}}}\xi^j\Sigma^+_{(i)}(X)\\
&=\coprod_{\substack{j\geq0\\0\leq k\leq i}}\xi^j\Sigma_k^+(X)
\end{aligned}\right\}\unsure{\underset{\mathrm{df}}{=}|\mathrel{=}|\xi^j}\Sigma^+_{[j]}(X)
\]
% archive sheet number: 38
